% !TeX root = ../../Book5.tex
% 纲要标题：### 20.2 中心作用
% 纲要位置：工作记录/计划纲要.md，第 4368 行。

\section{中心作用}
\label{b5:sec:2002}

% 纲要范围：
%
% * 无穷小特征标
% * 包络代数的中心
% * Harish–Chandra 同构的标准形式
%

Casimir 元在最高权模上给出一个标量，但不同最高权可能给出同一标量。中心的全部元素可以同时测量一个模，不过仍不能区分所有简单最高权模。本节将这种限制具体算出。

\subsection{无穷小特征标}
\label{b5:subsec:2002:01}

\begin{definition}\label{b5:def:infinitesimal-character}
设 \(\mathfrak g\) 是复李代数，\(Z=Z\Bigl(U(\mathfrak g)\Bigr)\)，\(M\ne0\) 是左 \(U(\mathfrak g)\) 模。若存在含幺代数同态 \(\chi:Z\to\mathbb C\)，使 \(zm=\chi(z)m\) 对所有 \(z\in Z,m\in M\) 成立，则称 \(M\) 具有\term[wuqiongxiao-tezhengbiao]{无穷小特征标}[infinitesimal character] \(\chi\)。这里所说的是包络代数中心的特征标，与群上取矩阵迹的特征标不同。
\end{definition}

\begin{proposition}\label{b5:prop:highest-central-character}
设 \(\mathfrak g\) 是有限维复半单李代数，固定 Cartan 子代数 \(\mathfrak h\) 和正根集合，\(\lambda\in\mathfrak h^*\)。每个最高权为 \(\lambda\) 的非零左 \(\mathfrak g\) 模 \(M\) 都有无穷小特征标，即包络代数中心 \(Z(U(\mathfrak g))\) 在 \(M\) 上按一个含幺代数同态 \(\chi_\lambda:Z(U(\mathfrak g))\to\mathbb C\) 作用。这个同态只依赖 \(\lambda\)，所有同一最高权的非零商模具有同一个 \(\chi_\lambda\)。
\end{proposition}
\begin{proof}
设 \(v\) 为最高权生成元。中心元 \(z\) 与 Cartan 子代数交换，所以保持一维空间 \(M_\lambda=\mathbb Cv\)，从而 \(zv=\chi_\lambda(z)v\)。对 \(u\in U(\mathfrak g)\)，有 \(zuv=uzv=\chi_\lambda(z)uv\)，故这个标量作用于全模。和、积及单位元在 \(v\) 上的作用说明 \(\chi_\lambda\) 是含幺代数同态。最后每个最高权模都是 \(M(\lambda)\) 的商，故标量与所选商无关。
\end{proof}

这个论证利用最高权空间的一维性及其生成全模的性质，无须将有限维 Schur 引理推广到无限维模。

\subsection{包络代数的中心}
\label{b5:subsec:2002:02}

固定三角分解。PBW 将包络代数写成有序负根、Cartan、正根单项式的线性组合。保留没有正负根因子的单项式，得到线性投影
\[
\operatorname{pr}:U(\mathfrak g)\longrightarrow U(\mathfrak h)
=\operatorname{Sym}(\mathfrak h).
\]
对中心元 \(z\)，它的伴随 Cartan 权为零。含正根因子的项零化最高权向量；只含非零负根因子的项不可能有权零。因此
\begin{equation}\label{b5:eq:central-character-projection}
\chi_\lambda(z)=\operatorname{pr}(z)(\lambda).
\end{equation}
投影 \(\operatorname{pr}\) 在整个 \(U(\mathfrak g)\) 上不是代数同态，但在中心上的乘法性由这个标量公式推出：两个多项式在全部 \(\lambda\in\mathfrak h^*\) 上相同，就确实相等。

\begin{theorem}\label{b5:thm:sl2-enveloping-center}
设 \(\mathfrak g=\mathfrak{sl}_2(\mathbb C)\)，取基 \(e,f,h\)，满足 \([h,e]=2e\)、\([h,f]=-2f\)、\([e,f]=h\)，并记 \(C\) 为 Killing 型对应的 Casimir 元。令
\[
\Omega=h^2+2h+4fe=8C.
\]
则 \(Z\Bigl(U(\mathfrak g)\Bigr)=\mathbb C[\Omega]\)。在最高权为 \(z\in\mathbb C\) 的模上，\(\Omega\) 作用为 \(z(z+2)\)。
\end{theorem}
\begin{proof}
\(\Omega=8C\) 的中心性来自\cref{b5:prop:casimir-central}和\eqref{b5:eq:sl2-killing-casimir}，作用标量由最高向量直接算得。现在证明没有遗漏其他中心元。

若 \(u\) 是非零中心元，取 PBW 滤过中的最高次主符号 \(p\in\mathbb C[e,f,h]\)。对每个 \(x\in\mathfrak g\)，\([x,u]=0\) 使 \(\operatorname{ad}x\) 零化 \(p\)。由 \(h\) 的权，\(p\) 中 \(e,f\) 次数相等，所以 \(p=F(t,h)\)，\(t=ef\)。再用 \(\operatorname{ad}e\)，得到
\[
0=e\left(h\frac{\partial F}{\partial t}-2\frac{\partial F}{\partial h}\right).
\]
多项式环无零因子，故括号为零。换元 \(q=h^2+4t\)，则
\(\mathbb C[t,h]=\mathbb C[q,h]\)，上述导子在新坐标中就是 \(-2\partial/\partial h\)。因此 \(p\in\mathbb C[q]\)。由于 \(p\) 齐次且 \(q\) 次数为二，必有 \(p=a q^r\)。\(\Omega\) 的主符号正是 \(q\)，于是 \(u-a\Omega^r\) 仍中心，且滤过次数严格下降。归纳至常数，得到 \(u\in\mathbb C[\Omega]\)。不同幂的主符号次数不同，也说明 \(\Omega\) 没有多项式关系。
\end{proof}

因此 \(\mathfrak{sl}_2\) 的两个最高权给出相同无穷小特征标，当且仅当
\[
z(z+2)=w(w+2),\qquad
w=z\ \text{或}\ w=-z-2.
\]
例如 \(L(0)\) 与无限维简单模 \(L(-2)=M(-2)\) 有同一中心特征标，却不同构。中心参数识别的是一个有限反射群轨道，而不是单个最高权。

\subsection{Harish–Chandra 同构的标准形式}
\label{b5:subsec:2002:03}

令 \(\rho=\frac12\sum_{\alpha>0}\alpha\)，在 \(\mathfrak h^*\) 上定义移位作用
\[
w\mathbin{\cdot}\lambda=w(\lambda+\rho)-\rho.
\]
这确为群作用，因为先平移 \(\rho\)、作通常 Weyl 作用、再平移回来。引入 \(\rho\) 是为了把刚才的 \(z\mapsto-z-2\) 化成关于原点的反射：\(z+1\mapsto-(z+1)\)。

\begin{theorem}[Harish--Chandra 同构]\label{b5:thm:harish-chandra}
设 \(\mathfrak g\) 是有限维复半单李代数，固定 Cartan 子代数 \(\mathfrak h\) 和正根集合 \(\Phi^+\)，\(W\) 为相应 Weyl 群，\(\rho=\frac12\sum_{\alpha\in\Phi^+}\alpha\)。记 \(\operatorname{pr}:U(\mathfrak g)\to U(\mathfrak h)=\operatorname{Sym}(\mathfrak h)\) 为 PBW 三角分解中保留纯 Cartan 项的线性投影，把 \(\operatorname{Sym}(\mathfrak h)\) 识别为 \(\mathfrak h^*\) 上的多项式函数，\(W\) 按其自然作用作用于这些函数。则映射
\[
\gamma:Z\Bigl(U(\mathfrak g)\Bigr)\longrightarrow
\operatorname{Sym}(\mathfrak h)^W,
\qquad
\gamma(z)(\mu)=\operatorname{pr}(z)(\mu-\rho)
\]
是代数同构。对 \(\lambda\in\mathfrak h^*\)，记 \(\chi_\lambda\) 为最高权 \(\lambda\) 的 Verma 模的无穷小特征标。则对任意中心元 \(z\)，有 \(\chi_\lambda(z)=\gamma(z)(\lambda+\rho)\)。对 \(\lambda,\nu\in\mathfrak h^*\)，\(\chi_\lambda=\chi_\nu\) 当且仅当 \(\nu=w\mathbin{\cdot}\lambda=w(\lambda+\rho)-\rho\) 对某个 \(w\in W\) 成立。
\end{theorem}

\begin{proofsketch}
前面已经由最高权向量的作用证明 \(\operatorname{pr}\) 在中心上是代数同态，且 \(\chi_\lambda(z)=\operatorname{pr}(z)(\lambda)\)。先核对平移后的 \(W\) 不变性。对支配整权 \(\lambda\) 及简单反射 \(s_\alpha\)，Verma 模 \(M(\lambda)\) 含有最高权为 \(s_\alpha\mathbin{\cdot}\lambda\) 的子模；同一个中心元在母模和子模上作用的标量相同，故
\[
 \operatorname{pr}(z)(\lambda)
 =\operatorname{pr}(z)(s_\alpha\mathbin{\cdot}\lambda).
\]
支配整权在 \(\mathfrak h^*\) 中 Zariski 稠密，而两端都是关于 \(\lambda\) 的多项式，所以等式对全部权成立。每个简单反射都满足这一式，故 \(\gamma(z)\) 属于 \(\operatorname{Sym}(\mathfrak h)^W\)。

再给包络代数和对称代数取次数滤过。PBW 对称化在每一层的主符号上是恒等，且与伴随 \(\mathfrak g\) 作用相容，因此中心的分次向量空间与 \(\operatorname{Sym}(\mathfrak g)^{\mathfrak g}\) 相应。若 \(z\) 的最高次符号是 \(p\)，投影到 Cartan 部分后，\(\gamma(z)\) 的最高次项正是 \(p\) 在 \(\mathfrak h^*\) 上的限制：PBW 重排产生的交换子以及平移 \(\rho\) 都只改变较低次数项。Chevalley 限制定理给出
\[
 \operatorname{Sym}(\mathfrak g)^{\mathfrak g}
 \xrightarrow{\ \sim\ }
\operatorname{Sym}(\mathfrak h)^W.
\]
这里可以用已经建立的有限维最高权表示说明限制同构。用 Killing 型识别李代数及其对偶。若一个伴随不变多项式在 \(\mathfrak h\) 上为零，它在 \(\mathfrak h\) 的全部幂零伴随指数像上也为零；\cref{b5:thm:cartan-conjugacy}证明中的开像计算说明这些像含有非空 Zariski 开集，因此原多项式为零，得到单射。

为证满射，记 \(P\subseteq\mathfrak h^*\) 为整权格。对每个次数 \(m\)，\(\operatorname{Sym}^m(\mathfrak h^*)\) 由 \(\nu^m\)、\(\nu\in P\)，线性生成：若一个线性泛函在这些幂上全为零，它在权坐标中的相应多项式就在整数格上恒为零，逐变量使用无限多个整数根便知该多项式为零。再对 \(W\) 平均，便知不变部分由轨道和 \(\sum_{\nu\in W\lambda}\nu^m\) 生成。

支配整权 \(\lambda\) 的简单模特征标具有三角展开
\[
 \operatorname{ch}L(\lambda)
 =\sum_{\nu\in W\lambda}e^\nu
 +\sum_{\mu<\lambda}a_{\lambda\mu}
   \sum_{\nu\in W\mu}e^\nu,
\]
其中 \(\mu\) 支配且 \(\lambda-\mu\) 是简单根的非负整数组合。低于固定 \(\lambda\) 的支配权有限：正定根内积给出 \(\|\mu\|\le\|\lambda\|\)，而权格离散。按这个有限偏序归纳，每个形式轨道和都是有限维模特征标的线性组合。在线性映射 \(e^\nu\mapsto\nu^m\) 下，特征标的像恰是
\(h\mapsto\operatorname{tr}(\rho(h)^m)\)。它延拓为 \(x\mapsto\operatorname{tr}(\rho(x)^m)\)，这是 \(\mathfrak g\) 上的伴随不变多项式，因为其沿 \([y,x]\) 的导数是交换子的迹，等于零。于是所有所需轨道和均能延拓，限制映射满射。

故 \(\gamma\) 的伴随分次映射是同构。对次数归纳，依次提升最高次不变多项式并减去已提升项，得到 \(\gamma\) 满射；若 \(\gamma(z)=0\)，其最高次符号为零，按次数下降得到 \(z=0\)，故它也单射。

最后，有限群的两个不同轨道可由不变多项式区分：先在两条有限轨道上作多项式插值，使取值分别为零和一，再对 \(W\) 平均。于是全部中心元在两个最高权模上取相同标量，恰当 \(\lambda+\rho\) 与 \(\nu+\rho\) 属于同一 \(W\) 轨道。这里省略了简单反射对应的 Verma 子模构造，以及 PBW 对称化与伴随作用相容的逐项展开；限制同构则由上面的稠密性、特征标三角性和迹多项式论证得到。限制论证可对照\cite[Section 23.1]{Humphreys1972LieAlgebras}，滤过比较和移位不变性见\cite[Section 8.8, Theorems 8.50, 8.52 and 8.54]{Kirillov2008LieGroups}。
\end{proofsketch}

前章 Weyl 特征标公式的证明不依赖这个一般同构。

\begin{example}\label{b5:exm:hc-sl3-orbit}
对 \(\mathfrak{sl}_3\)，把 \(\mathfrak h^*\) 写成 \(a+b+c=0\) 的三元组，\(W=S_3\) 排列三个坐标。其不变多项式由
\(a^2+b^2+c^2\) 与 \(a^3+b^3+c^3\) 生成：对称多项式基本定理把任一不变式写为初等对称式的多项式，而第一初等对称式为零，另外两个分别为 \(-\frac12(a^2+b^2+c^2)\)、\(\frac13(a^3+b^3+c^3)\)。二次 Casimir 只看到第一个参数，完整中心还需要三次参数。这里的取值点是 \(\lambda+\rho\)，不能漏掉平移。
\end{example}

无穷小特征标也提示了如何分组研究 \(\mathcal O\) 中的扩张。若中心并非按纯标量作用，可以要求每个 \(z-\chi(z)\) 在每个向量上局部幂零；这称为广义中心特征标条件。扩张可能保留广义特征标而失去纯标量作用，因此两个条件不能混写。
